\section{The control term: numerical check against the exact variance}\label{app:ctlnum}

A dense calculation checks Section~\ref{sec:controls}: 260 workers, 16 firms, 1{,}160 person-years, $k=2$ controls
correlated with the effects and varying within every match, Gaussian errors. The truth is \eqref{ctl:eq:exact},
computed exactly. Each estimator is averaged over 1{,}500 outcome draws with exact traces; unit variances are the
sample variance within every repeat-measurement set and the true variance elsewhere (the fitted model is not under
test). Entries of Table~\ref{tab:ctlnum} are the mean of the estimate over the exact variance. ``Oracle'' is
$4\hat m'S\hat m-2\tr(CSCS)$ with the true $S$ (a check of the calculation); ``known'' treats the control
coefficients as known.

\begin{table}[H]
\centering\small
\caption{Mean of the variance estimate over the exact variance of the fixed-offset estimate}\label{tab:ctlnum}
\begin{tabular}{@{}lllcccc@{}}
\toprule
Error variances & Units & Estimate & Worker & Firm & Covariance & Sum \\
\midrule
One per worker--firm cell & match & oracle & 1.000 & 0.996 & 0.992 & 1.000 \\
 & & known & 0.740 & 0.905 & 0.984 & 0.622 \\
 & & \eqref{ctl:eq:Vhat} & 0.999 & 0.994 & 0.992 & 0.998 \\
\addlinespace
 & observation & oracle & 1.000 & 1.000 & 1.001 & 1.000 \\
 & & known & 0.738 & 0.912 & 0.996 & 0.619 \\
 & & \eqref{ctl:eq:Vhat} & 1.002 & 1.004 & 1.005 & 1.002 \\
\addlinespace
Moving with a control & match & oracle & 0.997 & 0.994 & 0.989 & 0.997 \\
\quad inside the cell & & known & 0.684 & 0.847 & 0.983 & 0.537 \\
 & & \eqref{ctl:eq:Vhat} & 0.976 & 0.988 & 0.990 & 0.970 \\
\addlinespace
 & observation & oracle & 1.001 & 1.001 & 1.004 & 1.001 \\
 & & known & 0.676 & 0.853 & 1.002 & 0.528 \\
 & & \eqref{ctl:eq:Vhat} & 0.925 & 0.992 & 1.005 & 0.903 \\
\bottomrule
\end{tabular}
\end{table}

The shortfall in the last block is consistent with item~7 of Section~\ref{ctl:est}: under
observation deletion every worker--firm cell with two or more observations is a set, and here its variances move
with the control inside the cell; the control adjustment of the set variances and the unchanged pair term are also
active in that block. The bias of the point estimate, $\tr(CS)$, is 0.01 standard deviations. Under observation
units the jointly estimated version has 0.995 to 1.035 times the fixed-offset variance. The combined check uses
true variances outside sets and exact traces; it does not test the fitted model for units outside sets, and it
concerns the fixed-offset estimator, not the default jointly estimated one. Simulations through the command
(default joint estimation, fitted model, random draws) are summarised in Section~\ref{sec:impl-mc}.


\section{Designs, benchmarks and limits of the evidence}\label{app:evidence}

\subsection{Repeat measurements and the normal approximation}
The Veneto calibration fixes the estimated effects and assigns known noise variances.
The ten scenarios include a flat profile, quadratic and log-linear wage profiles,
fat-tailed and skewed errors, worker-specific variance factors, and profiles
that differ between movers and stayers. The two fitting steps and the cap on
estimated effects were developed on these designs. In one scenario the 1\%
of mover rows beyond the cap account for 33\% of the linear term of the
true variance of the estimated variance of the sum. Raising the cap from
three to six standard deviations changes the predicted standard-error ratio
at fixed characteristics from 0.95 to 0.99 in that scenario. The ratio
describes the model with known characteristics. Several additional changes
enter when outcomes and effects are estimated.

The six-year method study used exact stayer leverages, 1,600 fixed projections
for movers, and 16 draws for the trace. Controls and the centring constant
were treated as fixed. Its initial ten-scenario exercise set negative
eigenvalues of the estimated covariance matrix to zero. This can raise the
estimated variances. The fresh-draw validation omitted that step. It changed
the six-year ratios by at most 0.007, and the estimated matrix was positive
definite in all 4,400 validation draws across the three panels. The benchmarks
use projected weights and simulated traces at the true variances. For
non-Gaussian errors they omit cumulant terms arising from small nonzero
diagonal entries of the projected weights. They are approximate benchmarks
in those scenarios. Coverage is measured against the known component.

\subsection{A single bottleneck}
A separate study used two groups of 50 firms, with the number of workers
who moved between groups varying from 100 to 10, plus designs without a
bottleneck, skewed errors, unusually noisy bridging workers, and a stress
test with only three bridging workers. With 1,000 draws per design, the
prespecified $q=1$ intervals covered between 94.1\% and 100\% across the
75 registered component--design cases. In the ten-bridge designs with
normal errors in the main panel, the worst normal coverage was 48.7\%.
Comparing an interval based on estimated inputs with one based on the true
covariance showed that some of the extra length came from the geometry of
the interval itself.

The study missed its registered length tolerance: without a bottleneck, the
covariance interval in one design was 12.7\% longer than the normal interval,
above the 10\% limit. This is a test of the prespecified $q=1$ construction.
It did not test the automatic rule applied by the released package.
In the three-bridge stress test the condition that no single row carry the
leading score fails. The observed coverage in two Gaussian designs cannot
extend the coverage result to a fixed handful of bridging workers. The
simulations also did not stress the correlation boundary in
Corollary~\ref{cor:q1}. Two or more dominant directions lie outside the design.

\subsection{Package builds and comparability}
The automatic rule was chosen on 32,000 development fits. Each confirmation
manifest fixed its designs, seeds, deletion modes, repetitions and
acceptance thresholds in advance. The first no-control confirmation failed
on its result inventory when one fit returned point estimates only. The
second no-control confirmation used fresh designs, and the controls
confirmation used new seeds. A compatibility check carried the controls
result from \code{f3938986} to candidate \code{d4b2d942} on the basis of
36 byte-identical default results. The released plugins were later built
from \code{2082f209}. The check does not extend the registered confirmation
claim to the later compatibility and exact-eigensolver changes.
The \href{https://github.com/johannes-schmieder/fevc/tree/59bf4094/fevc/docs/default_interval}{package simulation records}
give the complete attempt counts and thresholds.
The \href{https://github.com/johannes-schmieder/fevc/blob/59bf4094/fevc/docs/repeat_measurement/unified/se_unified_20261116.pdf}{November 16 document}
gives the full scenario and component tables, including the alternative
covariance inputs and interval lengths.

\subsection{Interpreting the KSS comparison}
The rebuilt leave-one-out sample counts match KSS's published Table~I
to within 0.01\%. Random tie-breaking among the highest-paid spells can
change which firm a worker is assigned to. KSS do not publish corresponding
leave-two-out counts. The point estimates in the two-period levels and
first-difference formulations agree once the year-effect treatment and
numerical accuracy are aligned. At the default 200 projections, numerical
noise in the older-worker point estimates was about 10--12\% of their
standard errors. A comparison at 20,000 projections and with a fixed
offset brought the discrepancies within the numerical Monte Carlo error.
The exact two-period point estimate is invariant to centring here.
Centring can still change the raw products used for variance fitting and
hence the reported standard errors.

The simulation follows KSS's wage-change design with Student-$t_5$ errors,
scaled to the prescribed variances. It fits firm effects and the variance
equation on each older-worker sample and reruns the path searches on each
draw. KSS used a pooled fit and held the paths fixed. The primary design
puts all noise on the second year's wage. An alternative splits the noise
independently across the two years while preserving the variance of the
change. The first 100 leave-one-out and 200 leave-two-out draws were common
to all three procedures. The \fevc{} runs used 1,000 per sample. The
comparison counts an undefined interval, an interval with complex curvature
or an interval excluding its own point estimate as a miss under the scoring convention used in these comparisons. On the older leave-two-out sample, 19\%
of the KSS Matlab model-based $q=1$ intervals were scored invalid.
Coverage estimates based on 100--200 draws have Wilson half-widths of
roughly 2--6 percentage points in the reported cases.

The Econometrica replication code and the KSS Matlab package have different
implementation limits. The supplied critical-value table supports $q=1$
only, so the published $q=2$ intervals could not be reproduced. The
split-sample routine draws paths and traces without setting a seed, giving
differences of 0.5--8\% in standard errors across repeated calls on the
same data. Its remainder calculation indexes a single element where the
main routine uses the leading-direction column. The effect of this apparent
indexing error on interval widths has not been measured. In the replication
package's model-based routine, requesting the eigenvalue diagnostics
resulted in an error. The GitHub version of the KSS Matlab package supplied
the model-based $q=1$ intervals. Appendix~\ref{kss:app:repro} identifies
the actual programs, settings and inputs.
